% Part: normal-modal-logic
% Chapter: tableaux
% Section: soundness

\documentclass[../../../include/open-logic-section]{subfiles}

\begin{document}

\olfileid[fr]{nml}{tab}{sou}

\olsection{Correction pour \Log{K}}

\begin{explain}
Comme pour les tableaux propositionnels signés, on démontre la
correction en suivant une branche satisfaisable. Les règles
modales demandent aussi une interprétation des préfixes qui
respecte les relations entre un préfixe et ses successeurs.
\end{explain}

\begin{explain}
  Pour montrer que les !!{tableau}s préfixés sont corrects, nous devons
  montrer que, si
  \[
  \sFmla{\True}{!B_1}[1], \dots, \sFmla{\True}{!B_n}[1], \sFmla{\False}{!A}[1]
  \]
  admet un !!{tableau} fermé, alors $!B_1, \dots, !B_n \Entails !A$.
  Il est plus simple de démontrer la contraposée : s'il existe un
  $\mModel{M}$ et un monde~$w$ tels que $\mSat{M}{!B_i}[w]$ pour tout
  $i=1$, \dots,~$n$, mais $\mSat/{M}{!A}[w]$, aucun !!{tableau} ne peut
  se fermer. Ce contre-modèle montre que les hypothèses initiales du
  !!{tableau} sont satisfaisables. La stratégie consiste à montrer que,
  lorsque l’ensemble des !!{formula}s préfixées d'une branche d'%
  !!a{tableau} est satisfaisable, chaque application d'une règle
  produit au moins une branche prolongée qui l'est aussi. Comme les
  branches fermées ne sont pas satisfaisables, tout !!{tableau} issu
  d'un ensemble satisfaisable de !!{formula}s préfixées possède au
  moins une branche ouverte.

  Pour appliquer cette stratégie au cas modal, nous devons étendre la
  définition de « satisfaisable » aux formules modales munies de
  préfixes. Une fois cette définition posée, la démonstration est
  directe.
\end{explain}

\begin{defn}
  Soit $P$ un ensemble de préfixes, c'est-à-dire $P \subseteq (\PosInt)^*
  \setminus \{\emptyseq\}$, et soit $\mModel{M}$ un modèle. Une
  fonction~$f\colon P \to W$ est une \emph{interprétation de}~$P$
  dans~$\mModel{M}$ si, chaque fois que $\sigma$ et $\sigma.n$
  appartiennent tous deux à~$P$, on a $Rf(\sigma)f(\sigma.n)$.

  Relativement à une interprétation des préfixes de~$P$, nous posons :
  \begin{enumerate}
  \item $\mModel{M}$ satisfait $\sFmla{\True}{!A}[\sigma]$ si et seulement si
    $\mSat{M}{!A}[f(\sigma)]$.
  \item $\mModel{M}$ satisfait $\sFmla{\False}{!A}[\sigma]$ si et seulement si
    $\mSat/{M}{!A}[f(\sigma)]$.
  \end{enumerate}
\end{defn}

\begin{defn}
  Soit $\Gamma$ un ensemble de !!{formula}s préfixées, et soit $P(\Gamma)$
  l'ensemble des préfixes qui y figurent. Si $f$ est une interprétation
  de~$P(\Gamma)$ dans $\mModel{M}$, nous disons que $\mModel{M}$ satisfait
  $\Gamma$ relativement à~$f$, ce que nous notons $\mSat{M}{\Gamma}[f]$,
  si $\mModel{M}$ satisfait chaque !!{formula} préfixée de~$\Gamma$
  relativement à~$f$. L'ensemble $\Gamma$ est \emph{satisfaisable} s'il
  existe un modèle~$\mModel{M}$ et une interprétation~$f$ de
  $P(\Gamma)$ tels que $\mSat{M}{\Gamma}[f]$.
\end{defn}

\begin{prop}
  Si $\Gamma$ contient à la fois $\sFmla{\True}{!A}[\sigma]$ et
  $\sFmla{\False}{!A}[\sigma]$, pour une !!{formula}~$!A$ et un
  préfixe~$\sigma$, alors $\Gamma$ n'est pas satisfaisable.
\end{prop}

\begin{proof}
  Il ne peut exister un modèle~$\mModel{M}$ et une interprétation~$f$ de
  $P(\Gamma)$ tels que l'on ait à la fois $\mSat{M}{!A}[f(\sigma)]$ et
  $\mSat/{M}{!A}[f(\sigma)]$.
\end{proof}

\begin{thm}[Correction]
  \ollabel{thm:tableau-soundness}
  Supposons $P(\Gamma)$ clos par préfixes non vides. Si $\Gamma$ admet un !!{tableau} fermé,
  alors $\Gamma$ n'est pas satisfaisable.
\end{thm}

\begin{proof}
Nous disons qu'une branche d'!!a{tableau} est satisfaisable si et
seulement si l'ensemble des !!{signed formula}s qui y figurent est
satisfaisable ; !!a{tableau} est satisfaisable s'il possède au moins
une branche satisfaisable.

La clôture par préfixes non vides est préservée par les règles ; un préfixe nouveau n’a donc aucun descendant déjà présent. Montrons qu'une application d'une règle d'inférence à un !!{tableau}
satisfaisable produit toujours un !!{tableau} satisfaisable. Cela
établira le théorème : tout !!{tableau} fermé résulte d'applications
de règles d'inférence au !!{tableau} formé des seules hypothèses
de~$\Gamma$. Si $\Gamma$ était satisfaisable, tout !!{tableau}
construit à partir de lui le serait donc aussi. Or un !!{tableau}
fermé ne peut être satisfaisable : toutes ses branches sont fermées,
et une branche fermée n'est pas satisfaisable.

Supposons que nous ayons un !!{tableau} satisfaisable, c'est-à-dire
!!a{tableau} possédant au moins une branche satisfaisable. L'application d'une règle
d'inférence ajoute des !!{signed formula}s à une branche ou scinde
celle-ci en deux. Si le !!{tableau} possède une branche satisfaisable
qui n'est pas touchée par
cette application, elle reste satisfaisable dans le !!{tableau}
prolongé, qui est donc satisfaisable. Il suffit ainsi d'examiner le
cas où la règle s'applique à une branche satisfaisable.

Soit $\Gamma$ l'ensemble des !!{signed formula}s de cette branche, et
soit $\sFmla{S}{!A}[\sigma] \in \Gamma$ la !!{signed formula} à laquelle
s'applique la règle. Si la règle ne scinde pas la branche, nous devons
montrer que la branche prolongée, c'est-à-dire $\Gamma$ avec les
conclusions de la règle, reste satisfaisable. Si la règle scinde la
branche, nous devons montrer qu'au moins une des deux branches
obtenues est satisfaisable.

Examinons d'abord les règles qui ne scindent pas la branche.
\begin{enumerate}
\item On prolonge la branche en appliquant $\TRule{\True}{\lnot}$ à
  $\sFmla{\True}{\lnot !B}[\sigma] \in \Gamma$. La branche prolongée
  contient alors les !!{signed formula}s de $\Gamma \cup
  \{\sFmla{\False}{!B}[\sigma]\}$. Supposons $\mSat{M}{\Gamma}[f]$.
  En particulier, $\mSat{M}{\lnot !B}[f(\sigma)]$. Donc
  $\mSat/{M}{!B}[f(\sigma)]$, c'est-à-dire que $\mModel{M}$ satisfait
  $\sFmla{\False}{!B}[\sigma]$ relativement à~$f$.
\item On prolonge la branche en appliquant $\TRule{\False}{\lnot}$ à
  $\sFmla{\False}{\lnot !B}[\sigma] \in \Gamma$ : exercice.
\item On prolonge la branche en appliquant $\TRule{\True}{\land}$ à
  $\sFmla{\True}{!B \land !C}[\sigma] \in \Gamma$, ce qui ajoute deux
  !!{signed formula}s à la branche :
  $\sFmla{\True}{!B}[\sigma]$ et $\sFmla{\True}{!C}[\sigma]$. Supposons
  $\mSat{M}{\Gamma}[f]$, en particulier $\mSat{M}{!B \land
    !C}[f(\sigma)]$. Alors $\mSat{M}{!B}[f(\sigma)]$ et
  $\mSat{M}{!C}[f(\sigma)]$. Ainsi, $\mModel{M}$ satisfait à la fois
  $\sFmla{\True}{!B}[\sigma]$ et $\sFmla{\True}{!C}[\sigma]$
  relativement à~$f$.
\item On prolonge la branche en appliquant $\TRule{\False}{\lor}$ à
  $\sFmla{\False}{!B \lor !C}[\sigma] \in \Gamma$ : exercice.
\item On prolonge la branche en appliquant $\TRule{\False}{\lif}$ à
  $\sFmla{\False}{!B \lif !C}[\sigma] \in \Gamma$. Cela ajoute deux
  !!{signed formula}s à la branche : $\sFmla{\True}{!B}[\sigma]$ et
  $\sFmla{\False}{!C}[\sigma]$. Supposons
  $\mSat{M}{\Gamma}[f]$, en particulier
  $\mSat/{M}{!B \lif !C}[f(\sigma)]$. Alors
  $\mSat{M}{!B}[f(\sigma)]$ et
  $\mSat/{M}{!C}[f(\sigma)]$. Ainsi,
  $\mModel{M}, f$ satisfait à la fois
  $\sFmla{\True}{!B}[\sigma]$ et $\sFmla{\False}{!C}[\sigma]$.

\iftag{prvBox}{%
\item On prolonge la branche en appliquant $\TRule{\True}{\Box}$ à
  $\sFmla{\True}{\Box !B}[\sigma] \in \Gamma$ :
  \iftag{probBox}{Exercice.}{Cela ajoute une nouvelle !!{signed
      formula}~$\sFmla{\True}{!B}[\sigma.n]$ à la branche, pour un
    $\sigma.n \in P(\Gamma)$ (car $\sigma.n$ doit être déjà utilisé).
    Supposons $\mSat{M}{\Gamma}[f]$, en particulier $\mSat{M}{\Box
      !B}[f(\sigma)]$. Comme $f$ est une interprétation des préfixes
    et que $\sigma$, $\sigma.n \in P(\Gamma)$, nous savons que
    $Rf(\sigma)f(\sigma.n)$. Par conséquent,
    $\mSat{M}{!B}[f(\sigma.n)]$ ; autrement dit, $\mModel{M}, f$
    satisfait $\sFmla{\True}{!B}[\sigma.n]$.}

\item On prolonge la branche en appliquant $\TRule{\False}{\Box}$ à
  $\sFmla{\False}{\Box !B}[\sigma] \in \Gamma$ :
  \iftag{probBox}{Exercice.}{Cela ajoute une nouvelle !!{signed
      formula}~$\sFmla{\False}{!B}[\sigma.n]$, où $\sigma.n$ est un
    nouveau préfixe sur la branche, c'est-à-dire $\sigma.n \notin P(\Gamma)$.
    Puisque $\Gamma$ est satisfaisable, il existe un $\mModel{M}$ et
    une interprétation~$f$ de $P(\Gamma)$ tels que $\mSat{M}{\Gamma}[f]$,
    en particulier $\mSat/{M}{\Box !B}[f(\sigma)]$. Nous devons montrer
    que $\Gamma \cup \{\sFmla{\False}{!B}[\sigma.n]\}$ est
    satisfaisable. Pour cela, définissons une interprétation de
    $P(\Gamma) \cup \{\sigma.n\}$ comme suit :

  Puisque $\mSat/{M}{\Box !B}[f(\sigma)]$, il existe un $w \in W$ tel que
  $Rf(\sigma)w$ et $\mSat/{M}{!B}[w]$. Soit $f'$ identique à $f$, sauf
  que $f'(\sigma.n) = w$. Comme $f'(\sigma) = f(\sigma)$ et
  $Rf(\sigma)w$, nous avons $Rf'(\sigma)f'(\sigma.n)$, de sorte que $f'$
  est une interprétation de $P(\Gamma) \cup \{\sigma.n\}$. De plus,
  $\mSat/{M}{!B}[f'(\sigma.n)]$. Comme $f(\sigma') = f'(\sigma')$ pour
  tout préfixe $\sigma' \in P(\Gamma)$, $\mSat{M}{\Gamma}[f']$. Ainsi,
  $\mModel{M}, f'$ satisfait $\Gamma \cup
  \{\sFmla{\False}{!B}[\sigma.n]\}$.}
  }{}

\iftag{prvDiamond}{%
\item On prolonge la branche en appliquant $\TRule{\True}{\Diamond}$ à
  $\sFmla{\True}{\Diamond !B}[\sigma] \in \Gamma$ :
  \iftag{probDiamond}{Exercice.}{Cela ajoute une nouvelle !!{signed
      formula}~$\sFmla{\True}{!B}[\sigma.n]$, où $\sigma.n$ est un
    nouveau préfixe sur la branche, c'est-à-dire $\sigma.n \notin P(\Gamma)$.
    Puisque $\Gamma$ est satisfaisable, il existe un $\mModel{M}$ et
    une interprétation~$f$ de $P(\Gamma)$ tels que $\mSat{M}{\Gamma}[f]$,
    en particulier $\mSat{M}{\Diamond !B}[f(\sigma)]$. Nous devons montrer
    que $\Gamma \cup \{\sFmla{\True}{!B}[\sigma.n]\}$ est
    satisfaisable. Pour cela, définissons une interprétation de
    $P(\Gamma) \cup \{\sigma.n\}$ comme suit :

  Puisque $\mSat{M}{\Diamond !B}[f(\sigma)]$, il existe un $w \in W$ tel que
  $Rf(\sigma)w$ et $\mSat{M}{!B}[w]$. Soit $f'$ identique à $f$, sauf
  que $f'(\sigma.n) = w$. Comme $f'(\sigma) = f(\sigma)$ et
  $Rf(\sigma)w$, nous avons $Rf'(\sigma)f'(\sigma.n)$, de sorte que $f'$
  est une interprétation de $P(\Gamma) \cup \{\sigma.n\}$. De plus,
  $\mSat{M}{!B}[f'(\sigma.n)]$. Comme $f(\sigma') = f'(\sigma')$ pour
  tout préfixe $\sigma' \in P(\Gamma)$, $\mSat{M}{\Gamma}[f']$. Ainsi,
  $\mModel{M}, f'$ satisfait $\Gamma \cup
  \{\sFmla{\True}{!B}[\sigma.n]\}$.}
\item On prolonge la branche en appliquant $\TRule{\False}{\Diamond}$ à
  $\sFmla{\False}{\Diamond !B}[\sigma] \in \Gamma$ :
  \iftag{probDiamond}{Exercice.}{Cela ajoute une nouvelle !!{signed
      formula}~$\sFmla{\False}{!B}[\sigma.n]$ à la branche, pour un
    $\sigma.n \in P(\Gamma)$ (car $\sigma.n$ doit être déjà utilisé).
    Supposons $\mSat{M}{\Gamma}[f]$, en particulier
    $\mSat/{M}{\Diamond !B}[f(\sigma)]$. Comme $f$ est une interprétation
    des préfixes et que $\sigma$, $\sigma.n \in P(\Gamma)$, nous savons
    que $Rf(\sigma)f(\sigma.n)$. Par conséquent,
    $\mSat/{M}{!B}[f(\sigma.n)]$ ; autrement dit, $\mModel{M}, f$
    satisfait $\sFmla{\False}{!B}[\sigma.n]$.}
  }{}
\end{enumerate}
Examinons maintenant les règles qui scindent la branche en deux.
\begin{enumerate}
\item On prolonge la branche en appliquant $\TRule{\False}{\land}$ à
  $\sFmla{\False}{!B \land !C}[\sigma] \in \Gamma$, ce qui produit
  deux branches : la branche gauche passe par
  $\sFmla{\False}{!B}[\sigma]$ et la branche droite par
  $\sFmla{\False}{!C}[\sigma]$. Supposons $\mSat{M}{\Gamma}[f]$, en
  particulier $\mSat/{M}{!B \land !C}[f(\sigma)]$. Alors
  $\mSat/{M}{!B}[f(\sigma)]$ ou $\mSat/{M}{!C}[f(\sigma)]$. Dans le
  premier cas, $\mModel{M}, f$ satisfait $\sFmla{\False}{!B}[\sigma]$ ;
  la branche gauche est donc satisfaisable. Dans le second cas,
  $\mModel{M}, f$ satisfait $\sFmla{\False}{!C}[\sigma]$ ; la branche
  droite est donc satisfaisable.
\item On prolonge la branche en appliquant $\TRule{\True}{\lor}$ à
    $\sFmla{\True}{!B \lor !C}[\sigma] \in \Gamma$ : exercice.
\item On prolonge la branche en appliquant $\TRule{\True}{\lif}$ à
    $\sFmla{\True}{!B \lif !C}[\sigma] \in \Gamma$ : exercice.
\end{enumerate}
\iftag{prvFalse,prvTrue}{Les fermetures par constante sont impossibles à satisfaire par leur sémantique.}{} \iftag{prvIff}{Pour le biconditionnel vrai, les deux opérandes ont la même valeur : la branche prolongée avec leurs deux signes vrais, ou avec leurs deux signes faux, est satisfaite. Pour le biconditionnel faux, leurs valeurs sont opposées : l’une des deux branches prolongées avec ces signes est satisfaite. Le préfixe reste le même.}{}\end{proof}

\begin{prob}
Complétez la démonstration du \olref[nml][tab][sou]{thm:tableau-soundness}.
\end{prob}

\begin{cor}
\ollabel{cor:entailment-soundness}
Si $\Gamma \Proves !A$, alors $\Gamma \Entails !A$.
\end{cor}

\begin{proof}
  Si $\Gamma \Proves !A$, il existe des $!B_1$, \dots, $!B_n \in
  \Gamma$ tels que $\Delta = \{\sFmla{\False}{!A}[1], \sFmla{\True}{!B_1}[1],
  \dots, \sFmla{\True}{!B_n}[1]\}$ admet un !!{tableau} fermé. Nous
  voulons montrer que $\Gamma \Entails !A$. Supposons le contraire :
  pour un certain $\mModel{M}$ et un monde $w$, on a
  $\mSat{M}{!B_i}[w]$ pour $i=1$, \dots,~$n$, mais
  $\mSat/{M}{!A}[w]$. Posons $f(1) = w$ ; alors $f$ est une
  interprétation de~$P(\Delta)$ dans~$\mModel{M}$, et $\mModel{M}$
  satisfait~$\Delta$ relativement à~$f$. Or, d'après le
  \olref{thm:tableau-soundness}, $\Delta$ n'est pas satisfaisable,
  puisqu'il admet un !!{tableau} fermé : contradiction. Donc
  $\Gamma \Entails !A$.
\end{proof}

\begin{cor}
\ollabel{cor:weak-soundness}
Si $\Proves !A$, alors $!A$ est vraie dans tous les modèles.
\end{cor}
\begin{explain}
Si l'ensemble des préfixes est clos par segments initiaux non
vides, un préfixe nouveau ne peut avoir de descendant sur la
branche. On peut donc l'interpréter par un monde témoin accessible
depuis le monde de son parent, sans imposer de nouvelle
condition d'accessibilité vers un monde déjà attribué à un
préfixe de la branche.
\end{explain}
\end{document}
